\section*{Chapitre \ref{ch:rs:transaction-costs-in-research} --- Coûts de transaction dans la recherche}

\begin{solution}{exo:rs:transaction-costs-in-research:1}
$2 + 10\,000 \times 0.7 \times 0.02 \times \sqrt{0.05} = 2 + 31.3 = 33.3$ points de base.
\end{solution}

\begin{solution}{exo:rs:transaction-costs-in-research:2}
$0.491/(5.33 \times 252) = 3.66$ points de base par unité négociée~; sa demi-fourchette seule vaut 2, et l'impact à 1~milliard de dollars est très supérieur au reste.
\end{solution}

\begin{solution}{exo:rs:transaction-costs-in-research:3}
Séparément~: $2q^{3/2}$~; après compensation~: $(2q)^{3/2} = 2^{3/2}q^{3/2}$~; la transaction compensée coûte $(2^{3/2} - 2)/2 = 41~\%$ de plus.
\end{solution}

\begin{solution}{exo:rs:transaction-costs-in-research:4}
Son coût est dominé par le mouvement du cours pendant son exécution (63~points de base d'écart type à une participation de 1~\% contre 14 d'impact). Le coefficient est une
moyenne sur les ordres~: 500 ordres fixent $\eta$ à environ un tiers près (0.24 autour de 1.09), 5\,000 à environ 12~\% (0.08 autour de 0.71), 50\,000 à 4~\%.
\end{solution}

\begin{solution}{exo:rs:transaction-costs-in-research:5}
L'optimiseur négocie dès que la variation de la prévision couvre le coût. À 10~millions de dollars, les coûts sont faibles, de sorte qu'il suit le bruit quotidien de la prévision, qui n'est pas
de l'alpha, et paie 46.1~\% par an pour cela. À 1~milliard de dollars, les coûts le poussent à évoluer lentement, ce qui moyenne le bruit~; il gagne même davantage avant coûts.
\end{solution}

\begin{solution}{exo:rs:transaction-costs-in-research:6}
Pour $v > 0$, le minimiseur de $\frac12(x - v)^2 + s\lvert x\rvert + c\lvert x\rvert^{3/2}$ est positif ou nul~; si $v \le s$, il vaut 0~; sinon la condition du premier ordre
$x - v + s + \frac32 c\sqrt x = 0$ donne, avec $y = \sqrt x$, $y^2 + \frac32 cy - (v - s) = 0$, d'où $y = \frac12(-\frac32 c + \sqrt{\frac94 c^2 + 4(v - s)})$ et
$x = y^2$~; par symétrie, le signe de $v$ est rétabli.
\end{solution}

\begin{solution}{exo:rs:transaction-costs-in-research:7}
Coûts ignorés~: $-29.47$ brut, monte à 1.15 pour cinquante jours et descend à 0.69 pour deux cents jours. Coûts inclus~: 1.80 brut, 2.43 pour un jour, 0.94 pour deux cents jours.
L'optimiseur tenant compte des coûts lisse déjà le portefeuille par sa réticence à négocier, si bien qu'il n'a besoin que d'un peu de lissage de la prévision~; le portefeuille naïf tire tout son
lissage de la prévision.
\end{solution}

\begin{solution}{exo:rs:transaction-costs-in-research:8}
Le coût par unité négociée croît comme la racine carrée de la taille de chaque transaction, si bien que, en fraction du capital, les coûts sont multipliés par $\sqrt 2 = 1.41$ lorsque le fonds
double~; en dollars, les coûts sont multipliés par $2^{3/2} = 2.83$, tandis que les dollars bruts ne font que doubler. Le rendement du fonds après coûts baisse avec la taille, et au-delà d'une certaine taille son
gain en dollars baisse aussi (chapitre~28).
\end{solution}

\begin{solution}{pb:rs:transaction-costs-in-research:1}
\begin{enumerate}
\item Demi-fourchette de 2 points de base plus $0.7\sigma\sqrt{Q/V}$~; à 1~\% de participation et 2~\% de volatilité, 14 points de base d'impact contre un bruit de prix de 63.
\item 500~: $\eta = 1.09$ (0.24), exposant 0.59 (0.07)~; 5\,000~: 0.71 (0.08), 0.48 (0.03)~; 50\,000~: 0.67 (0.03), 0.49 (0.01).
\item Almgren et coauteurs~: une puissance 3/5 sur environ 700\,000 ordres Citigroup~; Tóth et coauteurs~: la racine carrée.
\item Les moyennes par classes éliminent l'essentiel du bruit de prix avant le passage aux logarithmes~; ce bruit croît avec la taille de l'ordre, si bien qu'une erreur type homoscédastique est
fausse.
\item Rotation de 5.33 par jour, coûts de 1\,223~\% par an, ratio de Sharpe net de $-29.47$.
\item $-0.31$, 0.06, 1.80 et 1.76.
\item Il négocie lentement et moyenne ainsi le bruit quotidien de la prévision.
\item Les coûts le conduisent à négocier 0.018 de son capital par jour~; son ratio de Sharpe avant coûts tombe à 2.11, le portefeuille ne pouvant suivre même les signaux lents, et le ratio net est de 1.76.
\item 8.27~\% et 8.73~\% du capital d'une équipe par an séparément~; 16.04~\% après compensation, soit une économie de 0.96 point (5.7~\%).
\item Le coût croît comme la puissance $3/2$ de la transaction~: la transaction combinée coûte plus cher que les deux transactions séparées (2.94 points par an ici).
\item Coûts ignorés~: $-29.47$, $-19.18$, $-10.62$, $-2.57$, 0.10, 1.02, 1.15, 0.93, 0.69 pour aucun lissage et des demi-vies de 1, 2, 5, 10, 20, 50, 100, 200 jours. Coûts
inclus~: 1.80, 2.43, 2.37, 2.17, 1.94, 1.71, 1.40, 1.15, 0.94.
\item \textbf{Résultat principal.} Le ratio de Sharpe net culmine à une demi-vie de cinquante jours (1.15) pour le portefeuille qui ignore les coûts et à un jour (2.43) pour le
portefeuille tenant compte des coûts, à 1~milliard de dollars.
\item Le portefeuille tenant compte des coûts, appliqué à la prévision lissée sur un jour (2.43).
\item Faible~: son ratio net passe de 1.78 à 10~millions de dollars, à 1.15 à 1~milliard de dollars et à $-0.23$ à 10~milliards de dollars.
\item Avec le modèle de coûts, la taille du fonds retenue, les coûts annuels, la rotation, le \omterm{def:rs:transaction-costs-in-research:smoothing}{coût de point mort} et les résultats nets à plusieurs tailles.
\item Lorsque les transactions se succèdent pendant le temps de décroissance de l'impact, comme c'est le cas du portefeuille quotidien~: un propagateur facture la seconde transaction pour
l'impact résiduel de la première transaction.
\item À partir des modèles de pré-négociation de ses courtiers ou d'estimations publiées, rapportées au volume et à la volatilité (invariance), et le remplacer par ses propres ajustements à mesure que les exécutions arrivent.
\item Leurs capacités ne sont pas additives~: les transactions dans le même sens s'additionnent selon une loi concave, si bien que la capacité de la société est inférieure à la somme.
\item Les deux ralentissent le portefeuille vers la composante persistante de la prévision~; le \omterm{def:rs:portfolio-construction-ii:aim}{portefeuille cible} le fait avec des taux de décroissance connus, le lissage s'opérant avec une demi-vie choisie.
\item Dans le modèle de marché, dans l'optimiseur et dans le signal, à la taille à laquelle la stratégie fonctionnera.
\end{enumerate}
\end{solution}

\begin{solution}{iq:rs:transaction-costs-in-research:1}
De manière concave~: l'impact moyen croît à peu près comme la racine carrée de la taille de l'ordre rapportée au volume quotidien, ajustée par la volatilité. Preuves~: Almgren et coauteurs sur
environ 700\,000 ordres Citigroup (une puissance 3/5 pour l'\omterm{def:rs:transaction-costs-in-research:impact}{impact temporaire}), Tóth et coauteurs (racine carrée), et études d'invariance des transitions de portefeuille.
\end{solution}

\begin{solution}{iq:rs:transaction-costs-in-research:2}
Mesurer sa rotation et son \omterm{def:rs:transaction-costs-in-research:smoothing}{coût de point mort}, le chiffrer à la taille du fonds à laquelle il tournera à l'aide d'un modèle d'impact ajusté, le construire en intégrant les coûts à l'optimiseur, le lisser,
et examiner la décroissance de son \omterm{def:rs:anatomy-of-a-predictor:ic}{coefficient d'information}~: un coefficient quotidien de 2.5 avant coûts peut devenir négatif après coûts.
\end{solution}

\begin{solution}{iq:rs:transaction-costs-in-research:3}
Enregistrer pour chaque ordre parent la décision et les \omterm{def:rs:simulation-versus-live:calibration}{prix d'arrivée}, les exécutions, la taille, la participation, la volatilité et le volume~; régresser le coût (net de la demi-fourchette) sur
$\sigma\sqrt{Q/V}$ avec des erreurs robustes et ajuster l'exposant sur des moyennes par classes~; contrôler le mouvement du marché~; utiliser des milliers d'ordres et réajuster régulièrement.
\end{solution}

\begin{solution}{iq:rs:transaction-costs-in-research:4}
Comme contrainte conique du second ordre ($t_i \ge \lvert d_i\rvert^{3/2}$ est représentable à l'aide de cônes tournés), ou, comme ici, par une méthode proximale~: le coût est séparable et
son pas proximal admet une forme close, un seuil doux suivi de la racine d'un trinôme du second degré en $\sqrt{\lvert d\rvert}$.
\end{solution}

\begin{solution}{iq:rs:transaction-costs-in-research:5}
Compenser leurs transactions en interne, répartir l'économie au prorata du coût autonome, et chiffrer chaque stratégie à son coût marginal pour la société~: les transactions dans le même sens coûtent plus cher
ensemble que séparément, si bien que la capacité combinée est inférieure à la somme.
\end{solution}

\begin{solution}{iq:rs:transaction-costs-in-research:6}
Si le rendement brut est $g$ par an sur le capital $A$ et que les coûts, en fraction du capital, croissent comme $k\sqrt A$ (les transactions croissent comme $A$, le coût par unité comme $\sqrt A$), le gain
en dollars vaut $gA - kA^{3/2}$, maximal pour $A^* = (2g/(3k))^2$, où le rendement du capital vaut $g/3$~; au-delà, chaque dollar ajouté réduit le bénéfice.
\end{solution}
